\documentclass[11pt]{article}
\usepackage[margin=0.85in]{geometry}
\usepackage[T1]{fontenc}
\usepackage{lmodern}
\usepackage{microtype}
\usepackage{booktabs}
\usepackage{tabularx}
\usepackage{xcolor}
\usepackage{hyperref}
\usepackage{siunitx}
\usepackage{enumitem}

\definecolor{accent}{HTML}{6A3FC7}
\definecolor{ink}{HTML}{202124}
\hypersetup{colorlinks=true,linkcolor=accent,urlcolor=accent}
\color{ink}
\sisetup{detect-all,group-separator={,},group-minimum-digits=4}
\setlist{nosep,leftmargin=*}
\setlength{\parskip}{0.55em}
\setlength{\parindent}{0pt}

\title{\textbf{wgfx Performance Relative to Raw WebGPU}\\
\large Order-Controlled Microbenchmarks and an End-to-End wDoom Case Study}
\author{e10b / wDoom project}
\date{10 October 2026}

\begin{document}
\maketitle

\begin{abstract}
This note evaluates whether the \texttt{wgfx} C++ convenience layer adds
material CPU overhead relative to direct WebGPU calls through
\texttt{wgpu-native}. In order-controlled microbenchmarks, the measured
differences for compute encoding and submission, storage updates, bind-group
creation, and render encoding are small, inconsistent in sign, or within the
observed run-to-run variability. For the operations tested, no repeatable
abstraction penalty is resolved: \texttt{wgfx} is effectively at performance
parity with raw WebGPU. A separate end-to-end wDoom experiment did identify a
real framework bottleneck---an unconditional completion wait after every
frame. Replacing it with bounded nonblocking submission nearly doubled warmed
throughput while preserving pixel-perfect output. This distinction is
important: the wrapper calls were not costly, but one synchronization policy
was.
\end{abstract}

\section{Claim}

The defensible conclusion is:

\begin{quote}
\textbf{For the operations and workloads measured, wgfx is effectively as fast
as raw WebGPU through wgpu-native. No repeatable core forwarding or encoding
overhead was resolved above the benchmark's noise floor.}
\end{quote}

This is not a claim that every possible program, resource-loading pattern, or
backend must behave identically. It is a measured statement about the tested
hot-path operations, under matched resources and workloads.

\section{Experimental setup}

\begin{itemize}
  \item Apple M3 Pro, 14 GPU cores, arm64; macOS 15.6.1.
  \item AppleClang 17.0.0; \texttt{RelWithDebInfo} build
        (\texttt{-O2 -g -DNDEBUG}).
  \item \texttt{wgpu-native} supplied by WebGPU-distribution commit
        \texttt{50010b083e4a1bbdac8f9102602df915fe7404cb}.
  \item Raw and wgfx paths share the same process, device, queue, shaders,
        buffers, textures, pipelines, bind-group shapes, and workload sizes.
  \item Each test performs ten warm-up batches followed by repeated timed
        samples using \texttt{std::chrono::steady\_clock}.
  \item Submission and render tests reverse raw/wgfx execution order on every
        sample. This prevents thermal, cache, and queue drift from consistently
        favoring either implementation.
  \item Five complete suite executions were collected. Values below average
        the per-run medians rather than treating all samples as independent.
\end{itemize}

Queue completion is performed outside CPU submission samples so neither path
inherits greater queue backpressure. The raw and wrapped paths ultimately issue
the same WebGPU operations; the comparison isolates wrapper-side work.

\section{Order-controlled results}

\begin{table}[ht]
\centering
\caption{CPU time per operation. Small positive and negative differences are
measurement variability, not evidence that either path intrinsically
accelerates WebGPU.}
\begin{tabularx}{\textwidth}{@{}X r r r@{}}
\toprule
Operation & Raw WebGPU & wgfx & Relative difference \\
\midrule
Compute encoding, one dispatch & \SI{5.763}{\micro\second} &
  \SI{5.779}{\micro\second} & $+0.3\%$ \\
Storage update, \SI{4}{\kibi\byte} & \SI{2.814}{\micro\second} &
  \SI{2.898}{\micro\second} & $+3.0\%$ \\
Bind-group recreation & \SI{1.767}{\micro\second} &
  \SI{1.708}{\micro\second} & $-3.3\%$ \\
Compute-pipeline creation & \SI{116.575}{\micro\second} &
  \SI{98.259}{\micro\second} & $-15.7\%$ \\
Encode + queue submit, one dispatch & \SI{16.625}{\micro\second} &
  \SI{16.359}{\micro\second} & $-1.6\%$ \\
Render encoding, one draw & \SI{12.192}{\micro\second} &
  \SI{11.983}{\micro\second} & $-1.7\%$ \\
Render encoding, 100 draws & \SI{22.033}{\micro\second} &
  \SI{22.550}{\micro\second} & $+2.3\%$ \\
Render encoding, 1,000 draws & \SI{127.008}{\micro\second} &
  \SI{129.558}{\micro\second} & $+2.0\%$ \\
\bottomrule
\end{tabularx}
\end{table}

The most suspicious original result was encode plus submission at
$+22.4\%$ for wgfx. That benchmark always ran raw first and wgfx second. After
reversing the order on every sample and increasing the sample count, the result
became $-1.6\%$. Across the five runs, its relative difference ranged from
$-3.8\%$ to $+1.3\%$. The sign changes, and the magnitude is small; therefore
there is no resolved submission overhead. The 1,000-draw render discrepancy
similarly contracted from an order-sensitive $-11.6\%$ to $+2.0\%$.

Pipeline creation's larger negative value is not interpreted as wgfx being
faster. Creation is relatively noisy, uses few samples, and is not a steady-state
frame-loop operation. It is evidence against claiming a stable penalty, not
evidence of acceleration.

\section{The bottleneck that was real}

The original \texttt{wgfx::frame()} implementation submitted and presented a
frame, then called both nonblocking and blocking device polling. Waiting for GPU
completion after every frame serialized CPU and GPU work. This was a policy
cost, not ordinary wrapper-call overhead.

An optional bounded nonblocking mode was added. It permits at most three
submitted frames, services completion callbacks without waiting when possible,
and waits only when the in-flight bound is full. Queue ordering makes reuse of
persistent buffers safe, and explicit draining protects reporting and shutdown.

In five alternating blocking/nonblocking wDoom demo-1 runs, the medians were:

\begin{table}[ht]
\centering
\begin{tabular}{@{}lrrr@{}}
\toprule
Metric & Blocking & Nonblocking & Change \\
\midrule
Throughput & 416.85 FPS & 786.62 FPS & $+88.7\%$ \\
Mean CPU frame time & \SI{2.399}{\milli\second} &
  \SI{1.271}{\milli\second} & $-47.0\%$ \\
Queue-completion latency & \SI{1.196}{\milli\second} &
  \SI{1.329}{\milli\second} & $+11.1\%$ \\
Maximum observed peak RSS & \SI{144.7}{\mebi\byte} &
  \SI{152.1}{\mebi\byte} & $+5.1\%$ \\
\bottomrule
\end{tabular}
\end{table}

The throughput improvement is large enough to identify the unconditional wait
as a genuine framework-level bottleneck. The modest latency and memory increases
are expected consequences of allowing several submissions to remain in flight.

\section{Correctness and end-to-end renderer comparison}

All four embedded Ultimate DOOM demos were replayed with nonblocking submission
and every-frame validation. The GPU result was compared against the deterministic
software-renderer reference:

\begin{itemize}
  \item 8,738 eligible frames completed;
  \item 559,232,000 indexed pixels compared;
  \item zero pixel mismatches;
  \item zero readback failures; and
  \item zero pending callbacks at completion.
\end{itemize}

With validation disabled, GPU walls completed the four-demo corpus in
\SI{8.660}{\second}, or 1,009.0 FPS aggregate. A matched run of the original CPU
wall loops, using the same Release binary and nonblocking presentation policy,
completed it in \SI{6.732}{\second}, or 1,297.9 FPS. Thus the synchronization
fix substantially narrowed the prior deficit, but the GPU wall renderer remains
22.3\% slower than the software wall loops in this workload. That remaining gap
does not contradict wrapper parity: it lies in renderer command preparation,
data transfers, and GPU execution rather than generic wgfx forwarding.

\section{Limitations}

\begin{itemize}
  \item Results apply to this machine, backend revision, compiler, and workload.
  \item CPU wall-clock timing cannot isolate shader-only GPU duration; GPU
        timestamp queries were unavailable in the requested feature set.
  \item Texture decoding, mip generation, unusual allocation patterns, and
        other APIs were not proven equivalent by these hot-path measurements.
  \item Differences of a few percent that change sign or move with test order
        should be treated as unresolved, not promoted into performance claims.
\end{itemize}

\section{Conclusion}

The evidence supports saying that \texttt{wgfx} is about as fast as direct
WebGPU for the core operations tested. Once test-order bias was removed, no
repeatable encoding or submission overhead remained. Profiling did uncover a
real synchronization mistake in the framework, and correcting it produced an
88.7\% median throughput gain without changing shaders, caches, rendering data,
or output pixels. That is stronger evidence than a synthetic parity claim
alone: the study distinguished negligible abstraction cost from a consequential
policy cost, fixed the latter, and verified both performance and correctness.

\section*{Reproducibility artifacts}

The repository contains the benchmark source, full methodology, and CSV data:

\begin{itemize}
  \item \texttt{deps/wgfx/benchmarks/wgfx\_benchmark.cpp}
  \item \texttt{deps/wgfx/PERFORMANCE.md}
  \item \texttt{deps/wgfx/benchmarks/results/m3-pro-order-controlled-summary.csv}
  \item \texttt{benchmarks/wgfx-frame-sync-2026-10-10.csv}
  \item \texttt{benchmarks/wgfx-four-demo-nonblocking-2026-10-10.csv}
\end{itemize}

\end{document}
